\documentclass[11pt, a4paper]{article}

% ── Packages ────────────────────────────────────────────────────────────────
\usepackage[margin=2.5cm]{geometry}
\usepackage{amsmath, amssymb}
\usepackage{booktabs}
\usepackage{xcolor}
\usepackage{parskip}
\usepackage{titlesec}
\usepackage{enumitem}
\usepackage{microtype}
\usepackage{graphicx}
\usepackage{listings}
\usepackage{array}

\definecolor{darkblue}{RGB}{0,50,100}
\definecolor{darkred}{RGB}{150,0,0}
\definecolor{mygreen}{RGB}{0,100,50}
\definecolor{lightblue}{RGB}{230,240,255}
\definecolor{lightyellow}{RGB}{255,252,230}
\definecolor{lightgreen}{RGB}{230,250,235}
\definecolor{lightred}{RGB}{255,235,235}
\definecolor{codegrey}{RGB}{245,245,245}

% ── Section formatting ───────────────────────────────────────────────────────
\titleformat{\section}{\large\bfseries\color{darkblue}}{}{0em}{}[\titlerule]
\titleformat{\subsection}{\normalsize\bfseries\color{darkblue}}{}{0em}{}

% ── Coloured left-bar boxes (reliable with list environments) ────────────────
\definecolor{gold}{RGB}{180,140,0}

\newenvironment{intubox}{%
  \vspace{4pt}\par\noindent%
  {\color{darkblue}\vrule width 3pt}\hspace{6pt}%
  \begin{minipage}{\dimexpr\linewidth-9pt\relax}\vspace{2pt}%
}{%
  \vspace{2pt}\end{minipage}\par\vspace{4pt}%
}

\newenvironment{notebox}{%
  \vspace{4pt}\par\noindent%
  {\color{gold}\vrule width 3pt}\hspace{6pt}%
  \begin{minipage}{\dimexpr\linewidth-9pt\relax}\vspace{2pt}%
}{%
  \vspace{2pt}\end{minipage}\par\vspace{4pt}%
}

\newenvironment{greenbox}{%
  \vspace{4pt}\par\noindent%
  {\color{mygreen}\vrule width 3pt}\hspace{6pt}%
  \begin{minipage}{\dimexpr\linewidth-9pt\relax}\vspace{2pt}%
}{%
  \vspace{2pt}\end{minipage}\par\vspace{4pt}%
}

\lstset{
  language=Python,
  basicstyle=\ttfamily\footnotesize,
  keywordstyle=\color{darkblue}\bfseries,
  commentstyle=\color{mygreen}\itshape,
  backgroundcolor=\color{codegrey},
  frame=single,
  framerule=0.4pt,
  breaklines=true,
  showstringspaces=false,
  numbers=none,
  xleftmargin=6pt,
  xrightmargin=6pt,
  aboveskip=4pt,
  belowskip=4pt,
}

\newcommand{\LSR}{\mathrm{LSR}}

% ════════════════════════════════════════════════════════════════════════════
\begin{document}

\begin{center}
  {\LARGE\bfseries\color{darkblue} The Lucas Critique in ML: Model Notes}\\[6pt]
  {\large An Intuitive Guide to Every Modelling and Estimation Choice}\\[4pt]
  {\normalsize 4815L \quad --- \quad D200 Machine Learning in Economics}
\end{center}

\medskip

\begin{notebox}
\textbf{Purpose of this document.}
This note walks through every modelling and estimation choice in the project
from first principles. The goal is intuition: why was this approach chosen,
what is the algorithm actually \emph{doing}, and how does it relate to the
others? Mathematics is kept to the minimum needed to be precise.
\end{notebox}

\tableofcontents

\newpage

% ════════════════════════════════════════════════════════════════════════════
\section{The Data-Generating Process (DGP)}
\label{sec:dgp}

\subsection{What it is}

The simulated data comes from a \textbf{Markov-switching AR(1)} process:
\[
  y_t = \mu_{s_t} + \phi_{s_t}\,y_{t-1} + \sigma_{s_t}\,\varepsilon_t,
  \qquad \varepsilon_t \overset{\text{iid}}{\sim} \mathcal{N}(0,1),
  \qquad s_t \in \{0,1\}
\]

The regime $s_t$ switches between 0 (expansion) and 1 (recession) according
to a Markov chain: the probability of staying in regime 0 is $P_{00}=0.95$ and
of staying in regime 1 is $P_{11}=0.90$. Think of $y_t$ as something like
monthly GDP growth: it follows an AR(1) in normal times, but with a different
mean, persistence, and volatility during recessions.

\begin{intubox}
\textbf{Intuition.}
The economy flips between two ``modes'' (expansion and recession). Within each
mode the series behaves like a standard autoregression, but the parameters ---
mean, persistence, volatility --- are different in each mode. The \emph{regime
identity} $s_t$ is \textbf{latent}: you can never directly observe whether you
are in regime 0 or 1. All you see is the time series $y_t$. Every model in this
project tries to infer $s_t$ from $y_t$ and use that inference to forecast.
\end{intubox}

\subsection{Why simulate?}

With simulated data the true $s_t$ is known. This lets us measure how accurately
each model \emph{recovers the regimes} (regime accuracy, Adjusted Rand Index),
not just how well it forecasts. We can also control the structural break exactly:
we know the true parameters before and after, so we know how hard the problem is.

\subsection{The structural break (Lucas shift)}

After the training period, the DGP parameters change. The regime labels stay the
same (still 0 and 1) but the mean, volatility, and transition probabilities shift.
Models trained on the old parameters must now forecast under the new ones without
being told anything changed. This is the Lucas Critique in action.

Two shift scenarios are used:

\begin{center}
\begin{tabular}{lcccc}
  \toprule
  Parameter & Pre-break & Mild post & Severe post & Effect \\
  \midrule
  $\mu_0$ (expansion mean) & 0.5 & 0.2 & $-$0.5 & Lower mean growth \\
  $\mu_1$ (recession mean) & $-$1.5 & $-$1.0 & $-$3.0 & Deeper recessions \\
  $\sigma_0$ (expansion vol) & 1.0 & 1.3 & 2.0 & More volatile \\
  $\sigma_1$ (recession vol) & 2.5 & 3.0 & 4.5 & Much more volatile \\
  $P_{00}$ (expansion persist.) & 0.95 & 0.90 & 0.80 & Shorter expansions \\
  \bottomrule
\end{tabular}
\end{center}

The ``mild'' shift resembles a gradual policy tightening. The ``severe'' shift
resembles the Volcker disinflation: a fundamental change in regime transition
dynamics on top of parameter changes.

% ════════════════════════════════════════════════════════════════════════════
\section{Evaluation Framework}

\subsection{The Lucas Sensitivity Ratio (LSR)}

All models are trained on pre-break data only, then frozen. The LSR compares
post-break forecast accuracy to pre-break accuracy:
\[
  \LSR = \frac{\text{RMSE}_\text{post}}{\text{RMSE}_\text{pre}}
\]

A model with LSR $>1$ performs worse after the break. LSR $=1$ means no
degradation. LSR $<1$ (which occurs for the Great Moderation break) means the
model actually \emph{improves} after the break because post-break data is less
noisy.

\begin{notebox}
\textbf{Important nuance.}
LSR is a \emph{relative} metric. A model with very low pre-break RMSE (because
it memorises training data) will mechanically have a high LSR, even if its
absolute post-break RMSE is reasonable. The XGB model is the clearest example:
pre-break RMSE of 0.40 produces LSR $=8.43$, even though the post-break RMSE of
3.34 is not the worst in absolute terms.
\end{notebox}

\subsection{Structural break tests}

Two formal tests are applied to detect whether a break occurred at all.

\textbf{Chow test.} Fits a single regression over the full sample and two
separate regressions (pre- and post-break). If the two sub-sample regressions
fit much better than the pooled one, the parameters must have shifted. The
$F$-statistic compares the reduction in residual sum of squares.

\begin{intubox}
\textbf{Intuition.}
Imagine fitting a straight line through two clouds of points and then fitting
two separate lines, one for each cloud. If the two-line fit is dramatically
better, the clouds must come from different underlying relationships. The Chow
test formalises this comparison.
\end{intubox}

\textbf{CUSUM test.} Computes the cumulative sum of prediction residuals over
time. Under parameter stability, the cumulative sum should wander randomly
around zero. A systematic drift away from zero signals that something has
changed. The test plots a running total and checks whether it strays outside
5\% significance bands.

The Chow test detects \emph{parameter} shifts in an AR model. It
does not detect a shift in \emph{transition dynamics} alone --- which is why it
correctly detects the Great Moderation (a variance shift) but fails to detect
the Volcker break (primarily a regime dynamics shift).

% ════════════════════════════════════════════════════════════════════════════
\section{Model 1: AR(2) and ARMA(2,1) Baselines}

\subsection{What they are}

The \textbf{AR(2)} model says the current value depends linearly on the two
most recent values plus noise:
\[
  y_t = c + \phi_1 y_{t-1} + \phi_2 y_{t-2} + \varepsilon_t
\]

The \textbf{ARMA(2,1)} adds a moving average term: the current error also
depends on last period's error:
\[
  y_t = c + \phi_1 y_{t-1} + \phi_2 y_{t-2} + \varepsilon_t + \theta\,\varepsilon_{t-1}
\]

\subsection{How they are estimated}

AR(2) is estimated by OLS: just regress $y_t$ on $y_{t-1}$ and $y_{t-2}$.
ARMA is estimated by maximum likelihood, using a Kalman filter to compute
the likelihood of the observed series given the parameters, then optimising
over $(\phi_1, \phi_2, \theta)$.

\subsection{Why include them?}

These are the \emph{benchmark}. They have no regime-switching at all. Their
purpose is to ask: does any of the complexity of regime-switching actually
help? If regime-switching models do not outperform a simple AR(2), the extra
complexity is not justified.

\begin{intubox}
\textbf{Intuition.}
AR(2) is essentially asking: ``given where I was in the last two periods, where
am I likely to be now?'' It treats the economy as a single, stable linear
process. The Lucas Critique predicts it will degrade after a structural break
--- and it does, with LSR around 1.66 for the mild shift. But because it fits
the pre-break data modestly (RMSE $\approx 1.63$), it does not have a large
denominator inflating its LSR.
\end{intubox}

% ════════════════════════════════════════════════════════════════════════════
\section{Model 2: Markov Switching Regression (MSM)}

\subsection{What it is}

The MSM is Hamilton's (1989) classical regime-switching model. It assumes the
data actually \emph{comes from} the Markov-switching DGP we described above:
\[
  y_t = \mu_{s_t} + \phi_{s_t}\,y_{t-1} + \sigma_{s_t}\,\varepsilon_t
\]
where $s_t$ follows a two-state Markov chain. This is the \emph{correctly
specified} model for our simulated data: the DGP is MS-AR(1), and so is the
MSM.

\subsection{How it is estimated: the Hamilton EM algorithm}

The challenge is that $s_t$ is unobserved. We cannot simply OLS-regress $y_t$
on $y_{t-1}$ because we do not know which regime's parameters to use at each
$t$.

The solution is the \textbf{EM (Expectation-Maximisation) algorithm}, a
general approach for estimating models with missing data. For the MSM:

\medskip
\textbf{E-step (Hamilton filter): infer regimes from data.}
Given current parameter guesses $(\hat\mu_k, \hat\phi_k, \hat\sigma_k, \hat{P})$,
compute the probability of being in each regime at each time step. This is done
recursively: the probability of being in state $k$ at time $t$ is proportional to

\begin{enumerate}[leftmargin=*, itemsep=2pt]
  \item How likely the observed $y_t$ is under regime $k$'s parameters
        (the Gaussian likelihood $\mathcal{N}(y_t;\,\hat\mu_k + \hat\phi_k y_{t-1},\,\hat\sigma_k^2)$).
  \item The probability of \emph{transitioning into} regime $k$ from wherever
        we were last period (the transition probability from the Markov chain).
\end{enumerate}

Running this forward through the data gives ``filtered'' regime probabilities.
Running a backward smoother then refines these using future data.

\medskip
\textbf{M-step: re-estimate parameters given inferred regimes.}
Using the regime probabilities from the E-step as soft weights, re-estimate
$(\mu_k, \phi_k, \sigma_k)$ by weighted least squares, and re-estimate $P$ by
counting weighted transition frequencies.

\medskip
Alternate E-step and M-step until convergence (parameters stop changing).

\begin{intubox}
\textbf{Intuition.}
Imagine trying to learn two different linear regression models from data where
you do not know which observations belong to which model. EM solves this by
iterating: first \emph{guess} which model each point belongs to (E-step), then
\emph{refit} the models using those guesses (M-step). The Hamilton filter is
the E-step for Markov-switching models: it propagates regime probabilities
forward through time, accounting for the Markov transition structure.

Think of the filter as a ``belief update'': at each time step, your prior belief
about the regime is updated by how well the new observation fits under each
regime's model. High volatility in $y_t$ raises the probability of being in the
high-variance recession regime.
\end{intubox}

\subsection{Why MSM has the worst Lucas performance despite correct specification}

This is one of the most striking findings. The MSM is \emph{exactly} the right
model for the simulated DGP, yet its post-break RMSE is the worst of all models
(5.97 mild, 11.96 severe).

The reason is the Hamilton filter. It uses the \emph{estimated} emission
densities $\mathcal{N}(y_t;\,\hat\mu_k + \hat\phi_k y_{t-1},\,\hat\sigma_k^2)$
to classify regimes. After the break, $\mu_k$ and $\sigma_k$ have changed, so
the filter's emission densities are wrong. The filter confidently assigns the
\emph{wrong} regime because it is looking for the old distribution. Post-break
regime accuracy collapses from 75\% to 20\% (effectively random). Once the
regime classification is wrong, the regime-specific forecasts become completely
wrong too.

\begin{notebox}
\textbf{Key lesson.}
Correct pre-break specification actually \emph{increases} Lucas vulnerability.
A model that commits strongly to the true pre-break structure has further to
fall when that structure changes. The MSM's precise identification of regimes
is a liability, not an asset, when the regime structure itself shifts.
\end{notebox}

% ════════════════════════════════════════════════════════════════════════════
\section{Model 3: Threshold Autoregression (TAR)}

\subsection{What it is}

The TAR model divides the data into two regimes based on whether the lagged
value of $y$ exceeds a threshold $\tau$:
\[
  y_t = \begin{cases}
    \hat\beta_0^\top x_t & \text{if } y_{t-1} \leq \hat\tau \\
    \hat\beta_1^\top x_t & \text{if } y_{t-1} > \hat\tau
  \end{cases}
\]

The threshold $\hat\tau$ is chosen by grid search: try every candidate
threshold (at quantiles from 15\% to 85\% of the distribution of $y_{t-1}$),
and pick the one that minimises the total sum of squared residuals across both
regimes. Features $x_t$ include $y_{t-1}$, rolling means, and covariates.

\subsection{How it is estimated}

\begin{enumerate}[leftmargin=*, itemsep=2pt]
  \item \textbf{Grid search for $\hat\tau$:} for each candidate threshold, split
        the data and fit OLS/Ridge in each partition. Compute total SSR. Pick the
        threshold that gives the lowest SSR.
  \item \textbf{Re-fit Ridge regressors} on each partition using the chosen $\hat\tau$.
\end{enumerate}

\begin{intubox}
\textbf{Intuition.}
The TAR says: ``the economy behaves differently when growth was recently weak
versus recently strong.'' The threshold is the boundary between these two
states. If $y_{t-1}$ was below $\hat\tau$ (bad times recently), use one
regression; if above (good times recently), use another.

The grid search is asking: what value of last period's output best separates
``expansion behaviour'' from ``recession behaviour''? It tries 50 candidate
split-points and picks the one that produces the cleanest partition.
\end{intubox}

\subsection{Why TAR is relatively Lucas-stable}

The threshold $\hat\tau$ is a hard cutoff in observable data ($y_{t-1}$), not
in a latent state that depends on fitted emission densities. After the break,
the TAR's regime classification is still driven by observable $y_{t-1}$
relative to a fixed threshold. The threshold may not perfectly separate
regimes any more, but it does not fail as catastrophically as the MSM's
Gaussian filter. LSR = 1.64 (mild), 2.78 (severe) --- among the best.

% ════════════════════════════════════════════════════════════════════════════
\section{Model 4: Hidden Markov Model + Ridge (HMM)}

\subsection{What it is}

The HMM is a \textbf{two-stage} pipeline:
\begin{enumerate}[leftmargin=*, itemsep=2pt]
  \item A \textbf{Gaussian HMM} infers latent regimes from an enriched
        feature vector.
  \item A \textbf{Ridge regressor} is fitted separately for each inferred regime.
\end{enumerate}

The HMM is estimated on the \emph{observation matrix}:
\[
  \mathbf{O}_t = \bigl[y_t,\;y_{t-1},\;\bar{y}^{(5)}_t,\;s^{(5)}_t,\;\bar{y}^{(20)}_t\bigr]^\top
\]
where $\bar{y}^{(5)}$ is the 5-period rolling mean, $s^{(5)}$ is the rolling
standard deviation, and $\bar{y}^{(20)}$ is the 20-period rolling mean.

\subsection{How it is estimated: Baum-Welch EM}

The HMM uses the same EM idea as the MSM, but the emission model is a
multivariate Gaussian over $\mathbf{O}_t$ rather than a univariate one over
$y_t$ alone.

\medskip
\textbf{E-step (Baum-Welch forward-backward algorithm):}
\begin{itemize}[itemsep=2pt]
  \item \emph{Forward pass}: compute $\alpha_t(k) = P(s_t=k,\,\mathbf{O}_{1:t})$
        recursively, using the multivariate Gaussian emission probability of $\mathbf{O}_t$.
  \item \emph{Backward pass}: compute $\beta_t(k) = P(\mathbf{O}_{t+1:T}\mid s_t=k)$.
  \item \emph{Smoothed posteriors}: $\gamma_t(k) = P(s_t=k\mid \mathbf{O}_{1:T}) \propto \alpha_t(k)\,\beta_t(k)$.
  \item \emph{Joint posteriors}: $\xi_t(j,k) = P(s_t=j,\,s_{t+1}=k\mid\mathbf{O}_{1:T})$.
\end{itemize}

\textbf{M-step:} update emission Gaussians (means and covariances) using the
soft regime weights $\gamma_t(k)$, and update the transition matrix from
$\xi_t(j,k)$.

After convergence, \textbf{Viterbi decoding} finds the single most probable
state sequence $s^*_{1:T}$, which is used to partition the training data for
fitting the Ridge regressors.

\textbf{At forecast time}, a \emph{causal forward filter} (no backward pass)
identifies the current regime from past observations only.

\subsection{The Viterbi algorithm}

Viterbi is \emph{dynamic programming} applied to hidden Markov models. It
finds the globally best path through the hidden states rather than the
best local state at each step.

\begin{enumerate}[leftmargin=*, itemsep=2pt]
  \item \textbf{Initialise:} $\log\delta_1(k) = \log\pi_k + \log b_k(\mathbf{O}_1)$.
  \item \textbf{Recurse:} $\log\delta_t(k) = \max_j[\log\delta_{t-1}(j) + \log A_{jk}] + \log b_k(\mathbf{O}_t)$.\\
        Store backpointer $\psi_t(k) = \arg\max_j[\log\delta_{t-1}(j) + \log A_{jk}]$.
  \item \textbf{Backtrack:} start at $s^*_T = \arg\max_k\log\delta_T(k)$, then
        follow backpointers backwards: $s^*_t = \psi_{t+1}(s^*_{t+1})$.
\end{enumerate}

Log-space is used throughout to avoid numerical underflow when multiplying
many small probabilities.

\begin{intubox}
\textbf{Intuition.}
Imagine a map where each column is a time step and each row is a possible
regime. Viterbi fills in the map column by column: the entry $\delta_t(k)$
records the probability of the best \emph{path} that ends at regime $k$ at
time $t$ (not just the local probability). When it reaches the final column,
it picks the best endpoint and traces back the entire path.

The key difference from the forward filter: Viterbi gives a \textbf{hard}
assignment (one definite regime per time step). The forward-backward smoother
gives \textbf{soft} probabilities ($\gamma_t(k)\in(0,1)$). Both are
$O(K^2 T)$ --- quadratic in regimes, linear in time.
\end{intubox}

\subsection{Why HMM is the most Lucas-stable model}

\begin{enumerate}[leftmargin=*, itemsep=3pt]
  \item \textbf{Richer observation matrix.} The HMM sees not just $y_t$ but
        also 5- and 20-period rolling statistics. These slow-moving signals
        carry information about the regime that does not change instantly after
        a parameter shift, giving the forward filter more stable inputs.
  \item \textbf{Soft regime assignments.} The forward filter outputs a
        probability distribution over regimes at each step, not a hard decision.
        After the break, when the emission densities are slightly wrong, the
        filter spreads probability mass across regimes rather than committing
        catastrophically to the wrong one.
  \item \textbf{Decoupled regime detection and forecasting.} The HMM identifies
        regimes independently of what the Ridge regressors predict. Even if the
        Ridge weights are slightly wrong post-break, the regime classification is
        not contaminated.
  \item \textbf{Ridge regularisation.} Ridge shrinks regression weights towards
        zero, reducing overfitting to the pre-break regime-specific patterns.
\end{enumerate}

% ════════════════════════════════════════════════════════════════════════════
\section{Model 5: Mixture of Experts (MoE)}

\subsection{What it is}

The Mixture of Experts is a \textbf{soft routing} model: rather than assigning
each observation to exactly one regime, it blends the predictions of $K$
``expert'' regressors using a \textbf{gating network} that assigns weights to
each expert:
\[
  \hat{y}_t = \sum_{k=1}^K \pi_k(x_t)\cdot\hat\beta_k^\top x_t
\]

The gating weights $\pi_k(x_t) = P(\text{expert}=k\mid x_t)$ come from a
logistic regression. The experts are Ridge regressors.

\subsection{How it is estimated: EM with soft responsibilities}

\textbf{E-step:} given current parameters, compute \emph{responsibilities}
(which expert should ``own'' each observation):
\[
  r_{ik} = \frac{\pi_k(x_i)\cdot\mathcal{N}(y_i;\,\hat\beta_k^\top x_i,\,\hat\sigma_k^2)}
              {\sum_j \pi_j(x_i)\cdot\mathcal{N}(y_i;\,\hat\beta_j^\top x_i,\,\hat\sigma_j^2)}
\]

The responsibility $r_{ik}$ is high for expert $k$ when: (a) the gating
network assigns high probability to $k$ at $x_i$, \emph{and} (b) expert $k$'s
prediction is close to the actual $y_i$.

\textbf{M-step:} re-estimate the gating network (logistic regression on
$(x_i, \arg\max_k r_{ik})$) and re-estimate each Ridge expert using the
responsibilities as sample weights.

\begin{intubox}
\textbf{Intuition.}
Think of three people (experts) each proposing a forecast, and a coordinator
(the gating network) deciding how much to trust each person based on the current
economic conditions. Each person learns only from the situations they are trusted
to handle. The EM algorithm refines both the experts and the coordinator
simultaneously: if an expert is more accurate, it gets more responsibility,
which makes its weights better, which makes it more accurate.

Unlike TAR or MSM, \emph{every} expert contributes to every forecast
proportionally to its gating weight. This ``soft'' blending provides partial
regularisation: a misclassified observation in the wrong regime does not
completely corrupt the forecast, it just reduces the correct expert's weight.
\end{intubox}

\subsection{Difference from HMM}

The MoE has no Markov chain. Each time step's regime probability is determined
solely by the current $x_t$, with no memory of the previous regime. The HMM
tracks how regimes evolve over time via the transition matrix; the MoE treats
each step independently. This makes MoE faster to estimate but means it cannot
learn that ``recessions last 10 periods on average.''

% ════════════════════════════════════════════════════════════════════════════
\section{Model 6: ML Regime (XGB)}

\subsection{What it is}

The ML Regime model uses a three-stage pipeline:
\begin{enumerate}[leftmargin=*, itemsep=3pt]
  \item \textbf{KMeans clustering}: cluster training observations into
        $K=2$ groups using lags, rolling statistics, and covariates. These
        clusters serve as pseudo-regime labels. No economics assumed ---
        purely data-driven.
  \item \textbf{XGBClassifier}: train a gradient-boosted classifier to
        reproduce those cluster labels from the same feature set. This gives
        a smooth, interpretable mapping from features to regime probabilities.
  \item \textbf{Per-regime XGBRegressors}: train one XGBoost regressor on
        each regime's subset of training data.
\end{enumerate}

\subsection{How XGBoost works}

XGBoost builds an \textbf{ensemble of decision trees} in sequence. Each tree
tries to correct the errors of all previous trees. The prediction is the sum of
all tree outputs:
\[
  \hat{y}_t = \sum_{m=1}^M f_m(x_t)
\]
where each $f_m$ is a small decision tree. Trees are added greedily: at each
step, find the tree $f_m$ that most reduces the loss on the current residuals.

Decision trees work by splitting the feature space: ``if rolling std $> 2.1$
AND lag-1 $< -0.5$, predict $-1.3$; else predict $+0.4$''. XGBoost's output is
therefore a \textbf{piecewise-constant function}: the feature space is divided
into rectangular regions, each with a constant prediction.

\begin{intubox}
\textbf{Intuition.}
Imagine drawing a very detailed decision flowchart for predicting output:
``If the economy was below $-0.5$ last month AND volatility was above $2.1$,
AND it rained in Des Moines...'' With 200 trees and no restraint, XGBoost can
carve the feature space into tiny boxes, each perfectly fitting the training
observations inside it. This is the memorisation that causes its spectacular
collapse after the structural break.

The piecewise-constant nature is the key vulnerability. Once the DGP shifts,
the box boundaries (learned from pre-break data) no longer correspond to
meaningful divisions. But because the boxes are rigid, the model cannot
interpolate between them --- it just outputs the wrong pre-break average for
each new observation that lands in an old box.
\end{intubox}

\subsection{Why XGB is the most Lucas-vulnerable model}

\begin{enumerate}[leftmargin=*, itemsep=3pt]
  \item \textbf{Frozen KMeans boundaries.} The regime classifier uses cluster
        centroids fixed at training time. Post-break, the same boundaries
        systematically misclassify regimes.
  \item \textbf{Piecewise-constant extrapolation.} Unlike polynomials or neural
        networks, tree predictions do not interpolate smoothly between training
        points. A post-break observation that falls in a region with no training
        data gets the pre-break average for that region --- which may be
        arbitrarily wrong.
  \item \textbf{Extreme in-sample fit.} Pre-break RMSE of 0.40 --- about $4\times$
        lower than the next best model --- reflects memorisation of the training
        data. This also inflates the LSR denominator, making the ratio look worse
        than it is in absolute terms.
\end{enumerate}

% ════════════════════════════════════════════════════════════════════════════
\section{Model 7: Markov Switching Neural Network (MSNN)}

\subsection{What it is}

The MSNN combines:
\begin{itemize}[itemsep=3pt]
  \item \textbf{A Markov chain} governing regime transitions (like the MSM).
  \item \textbf{Per-regime MLP experts} that produce forecasts (instead of
        linear regressions).
\end{itemize}

The data-generating assumption is:
\[
  y_t = f_{s_t}(x_t) + \sigma_{s_t}\,\varepsilon_t, \qquad s_{t}\mid s_{t-1}\sim P
\]
where each $f_k$ is a small feedforward neural network.

\subsection{The MLP (Multi-Layer Perceptron)}

Each expert is an MLP with architecture $n_\text{feat} \to 32 \to 16 \to 1$:

\begin{itemize}[itemsep=2pt]
  \item \textbf{Input layer}: the feature vector $x_t$ (lags, rolling stats, covariates).
  \item \textbf{Hidden layer 1 (32 neurons)}: each neuron computes $\text{ReLU}(w^\top x + b)$.
        ReLU means ``zero if negative, identity if positive'' --- it introduces
        nonlinearity while remaining computationally simple.
  \item \textbf{Hidden layer 2 (16 neurons)}: same structure.
  \item \textbf{Output layer (1 neuron)}: linear combination of the 16 hidden
        activations, producing a scalar forecast.
\end{itemize}

\begin{intubox}
\textbf{Intuition.}
The MLP can represent arbitrarily complex nonlinear functions, given enough
neurons. Think of it as a very flexible curve-fitter: rather than assuming
$y_t$ depends \emph{linearly} on lagged values, the MLP can learn any smooth
relationship from the training data. Mathematically, a two-layer ReLU network
is a piecewise-\emph{linear} function of the inputs (unlike XGBoost's
piecewise-\emph{constant} output), which allows it to interpolate between
training points rather than jumping to flat plateaus.
\end{intubox}

\subsection{How it is estimated: EM with Adam}

The MSNN uses the same EM structure as the MSM and HMM, but the M-step updates
MLPs via gradient descent instead of closed-form weighted OLS.

\medskip
\textbf{E-step (Hamilton forward-backward filter):}
Same as the MSM, but the emission probability $P(y_t\mid s_t=k)$ is computed
from the MLP's prediction: $\mathcal{N}(y_t;\,f_k(x_t),\,\sigma_k^2)$.

Run the forward filter:
\[
  \log\alpha_t(k) = \log\sum_j e^{\log\alpha_{t-1}(j) + \log P_{jk}} + \log p(y_t\mid s_t=k)
\]
Then the backward filter, and combine to get smoothed posteriors $\gamma_t(k)$.

\textbf{M-step (Adam for MLP, analytical for transition matrix):}

\begin{itemize}[itemsep=3pt]
  \item \textbf{Transition matrix}: $\hat{P}_{jk} \propto \sum_t \xi_t(j,k)$.
        This is just counting weighted transitions --- no gradient needed.
  \item \textbf{MLP experts}: minimise the weighted loss
        $\mathcal{L}_k = \sum_t \gamma_t(k)\,(y_t - f_k(x_t))^2$
        using the \textbf{Adam} optimiser. Adam adjusts the learning rate
        individually for each weight based on the recent history of gradients.
\end{itemize}

\begin{intubox}
\textbf{Adam intuition.}
Imagine walking downhill (gradient descent) to find the lowest point in a
landscape. Vanilla gradient descent takes equal-sized steps in the direction of
steepest descent. Adam keeps a running average of the gradient direction
(\emph{momentum}) and a running average of the squared gradient magnitude
(\emph{adaptive scaling}). This means it takes bigger steps when the gradient
has been consistently pointing in one direction, and smaller steps when
gradients have been noisy. In practice this converges much faster and more
reliably than plain gradient descent for neural networks.
\end{intubox}

\subsection{The MSNN's paradox}

The MSNN has:
\begin{itemize}[itemsep=2pt]
  \item \textbf{Worst regime identification}: in-sample regime accuracy $\approx 51\%$
        (near random). The EM collapsed to a single dominant regime; the Markov
        chain stopped being meaningful.
  \item \textbf{Best absolute post-break RMSE}: 2.493 (mild), 3.391 (severe) --- the
        lowest of all models in every experiment.
\end{itemize}

How? Because the MLP experts are smooth universal approximators. Even if the
Markov structure has collapsed, the MLP can learn a good overall nonlinear
mapping from $x_t$ to $y_t$. Post-break, the MLP \emph{extrapolates smoothly}
outside the training distribution, whereas XGBoost's step functions return
stale pre-break box averages.

\begin{notebox}
\textbf{Why smooth extrapolation is important for the Lucas Critique.}
After a structural break, test observations lie in a region of feature space
that training data only partially covers. Smooth functions (MLPs) extrapolate
by continuing the learned smooth relationship into this new region. Piecewise-
constant functions (XGBoost trees) return the average of the nearest training
box, which may be far from the new distribution. This is why the MSNN wins on
absolute accuracy despite its poor in-sample regime identification.
\end{notebox}

% ════════════════════════════════════════════════════════════════════════════
\section{Model 8: Model Average (Ensemble)}

\subsection{What it is}

The model average simply takes the \textbf{equal-weight mean} of all other
models' predictions at each time step:
\[
  \hat{y}_t^\text{avg} = \frac{1}{N}\sum_{n=1}^N \hat{y}_t^{(n)}
\]

No additional estimation is needed --- it is computed at prediction time only.

\subsection{Why it works}

The theoretical justification comes from the bias-variance trade-off. If
individual models have correlated errors, averaging reduces variance but does
not hurt bias much. But more relevantly here: the forecast risk decomposition
shows that the model class $\mathcal{F}$ contributes an \emph{approximation
error} term. Averaging across models with different functional forms effectively
expands $\mathcal{F}$, reducing the approximation error component.

\begin{intubox}
\textbf{Intuition.}
If one expert is wrong in one direction and another is wrong in the other
direction, averaging them gets you closer to the truth. More formally: even if
no single model has the right functional form for the new regime, the average
of models with diverse functional forms is likely to be closer to the truth
than any individual model. In 2 out of 4 experiments, the ensemble produces
both a lower LSR and lower absolute post-break RMSE than most individual models.
\end{intubox}

% ════════════════════════════════════════════════════════════════════════════
\section{Feature Engineering}

All models (except AR/ARMA and the raw MSM) use the same feature set:

\begin{center}
\begin{tabular}{lll}
  \toprule
  Feature & Formula & Intuition \\
  \midrule
  \texttt{y\_lag1} & $y_{t-1}$ & Most recent value \\
  \texttt{y\_lag2} & $y_{t-2}$ & Two-period lag \\
  \texttt{roll\_mean\_5} & $\bar{y}^{(5)}_t = \frac{1}{5}\sum_{i=0}^4 y_{t-i}$ & Short-run trend \\
  \texttt{roll\_std\_5} & $s^{(5)}_t$ & Short-run volatility \\
  \texttt{roll\_mean\_20} & $\bar{y}^{(20)}_t$ & Long-run trend \\
  \texttt{roll\_std\_20} & $s^{(20)}_t$ & Long-run volatility \\
  \texttt{exog\_0} & EFFR (simulated) / EFFR (real) & Interest rate level \\
  \texttt{exog\_1} & Simulated / UNRATE & Labour market \\
  \bottomrule
\end{tabular}
\end{center}

The rolling statistics are particularly important for the HMM. The HMM
observation matrix uses $[\bar{y}^{(5)}, s^{(5)}, \bar{y}^{(20)}]$ rather
than raw $y$ values, because regime transitions are more reliably detected from
slow-moving summary statistics than from a single noisy observation.

% ════════════════════════════════════════════════════════════════════════════
\section{Summary: How Each Model Answers the Fundamental Question}

\begin{center}
\small
\begin{tabular}{p{2.5cm}p{3.5cm}p{3.5cm}p{3.5cm}}
  \toprule
  \textbf{Model} & \textbf{How does it find regimes?} & \textbf{How does it forecast?} & \textbf{Lucas vulnerability} \\
  \midrule
  AR(2) / ARMA & Doesn't --- one global model & Linear in lags & Moderate: no regime assumptions to break \\
  \addlinespace
  MSM & Hamilton EM filter on $y_t$ & Regime-probability-weighted AR & \textcolor{darkred}{High}: emission densities wrong post-break; filter collapses \\
  \addlinespace
  TAR & Grid-search threshold in $y_{t-1}$ & Per-regime Ridge (hard split) & Low-moderate: observable threshold, no latent state to misidentify \\
  \addlinespace
  HMM & Baum-Welch EM on rolling-feature matrix & Per-regime Ridge (soft split) & \textcolor{mygreen}{Low}: soft posteriors, rich observation matrix \\
  \addlinespace
  MoE & Logistic gating via EM & Responsibility-weighted Ridge & Moderate: soft but no Markov memory \\
  \addlinespace
  ML XGB & KMeans $\to$ XGB classifier & Per-regime XGBRegressor & \textcolor{darkred}{Very high}: frozen boundaries + piecewise constant \\
  \addlinespace
  MSNN & Markov + MLP (EM) & Per-regime MLP & High LSR / \textcolor{mygreen}{Low abs. RMSE}: MLP smooth extrapolation \\
  \addlinespace
  Ensemble & --- & Equal-weight average & Low: diversification reduces approximation error \\
  \bottomrule
\end{tabular}
\end{center}

\bigskip

\begin{greenbox}
\textbf{The big picture.}
Every model in this project answers the same question: ``given what I know
about the economy's past, which \emph{mode} is it currently in, and how should
I forecast from that mode?''

The models differ in three key dimensions:
\begin{enumerate}[itemsep=2pt]
  \item \textbf{How they find the mode} (hard threshold, probabilistic filter,
        clustering, EM with Markov chain).
  \item \textbf{How flexible the forecast is} within each mode (OLS, Ridge,
        XGBoost trees, MLP).
  \item \textbf{How much they commit to the pre-break mode structure} (MSM
        commits maximally; HMM commits softly; MSNN commits to functional form
        but not regime structure).
\end{enumerate}

The central finding --- flexible latent-state estimation beats flexible
functional form --- can be rephrased simply: it is better to remain uncertain
about \emph{which mode you are in} (HMM's soft posteriors) than to be very
certain about \emph{what function to use} within a mode (XGBoost's trees).
After a structural break, the mode boundaries shift. Uncertain regime
assignments partially absorb this; rigid functional forms cannot.
\end{greenbox}

\end{document}
